\documentclass[a4paper,11pt]{article}

\usepackage[ngerman]{babel}
% Shared preamble for all CV languages.
% Load babel with the right language in each cv.tex BEFORE \input-ing this file.

\usepackage{latexsym}
\usepackage[empty]{fullpage}
\usepackage{titlesec}
\usepackage{marvosym}
\usepackage[usenames,dvipsnames]{color}
\usepackage{verbatim}
\usepackage{enumitem}
\usepackage[hidelinks]{hyperref}
\usepackage{fancyhdr}
\usepackage{tabularx}
\usepackage{multicol}
\input{glyphtounicode}

\usepackage[sfdefault]{inter}
\usepackage[T1]{fontenc}


\pagestyle{fancy}
\fancyhf{}
\fancyfoot{}
\setlength{\footskip}{10pt}
\renewcommand{\headrulewidth}{0pt}
\renewcommand{\footrulewidth}{0pt}

% Top margin is set per language in each cv.tex
\addtolength{\oddsidemargin}{0.0in}
\addtolength{\evensidemargin}{0.0in}
\addtolength{\textwidth}{0.0in}
\addtolength{\textheight}{0.0in}


\urlstyle{same}

%\raggedbottom
\raggedright
\setlength{\tabcolsep}{0in}

\titleformat{\section}{
  \it\vspace{3pt}
}{}{0em}{}[\color{black}\titlerule\vspace{-5pt}]

\pdfgentounicode=1

\newcommand{\resumeItem}[1]{
  \item{
    {#1\vspace{-4pt}}
  }
}

\newcommand{\resumeSubheading}[4]{
  \vspace{-2pt}\item
    \begin{tabular*}{0.97\textwidth}[t]{l@{\extracolsep{\fill}}r}
      \textbf{#1} & #2 \\
      \textit{\small #3} & \textit{\small #4} \\
    \end{tabular*}\vspace{-10pt}
}


\newcommand{\resumeSubItem}[1]{\resumeItem{#1}\vspace{-3pt}}
\renewcommand\labelitemii{$\vcenter{\hbox{\tiny$\bullet$}}$}
\newcommand{\resumeSubHeadingListStart}{\begin{itemize}[leftmargin=0.15in, label={}]}
\newcommand{\resumeSubHeadingListEnd}{\end{itemize}}
\newcommand{\resumeItemListStart}{\begin{itemize}}
\newcommand{\resumeItemListEnd}{\end{itemize}\vspace{-2pt}}


\addtolength{\topmargin}{-0.6in}  % pull content up
\addtolength{\textheight}{0.3in}  % use some of the (empty) bottom margin

\begin{document}


\begin{center}
    { \textbf{\Huge Mert Önder}} \\ \vspace{3pt}
   \small +49 17684492334 $|$ \underline{\href{mailto:mertonder51@gmail.com}{\nolinkurl{mertonder51@gmail.com}}} \\ \underline{\href{https://www.linkedin.com/in/mert--onder/}{\nolinkurl{linkedin.com/in/mert--onder}}} $|$ \underline{\href{https://github.com/MertOender}{\nolinkurl{github.com/MertOender}}}
\end{center}


%-----------BILDUNGSWEG-----------
\section{Bildungsweg}
\resumeSubHeadingListStart
    \resumeSubheading
        {Technische Universität München}{April 2026 -- heute}
        {M.Sc. Robotics, Cognition, Intelligence}{München, Deutschland}
       \vspace{0.1cm}
    \resumeSubheading
        {Technische Universität München}{Oktober 2022 -- April 2026}
        {B.Sc. Elektrotechnik und Informationstechnik}{München, Deutschland}
        \resumeItemListStart
       \small\resumeItem{Abschlussnote: 1,9 (beste 15\,\% des Jahrgangs)}
       \vspace{0.05cm}
       \resumeItem{Bachelorarbeit: \textit{A Modular ROS2 Architecture for SSVEP Based Brain Computer Interface to Balance Timing, Efficiency and Observability} (1,0)}
       \resumeItemListEnd
       \vspace{0.1cm}
    \resumeSubheading
        {Deutsche Schule Istanbul}{September 2017 -- Juni 2022}
        {Deutsche Auslandsschule - Abitur: 1,5}{Istanbul, Türkei}
        \vspace{0.1cm}
    \resumeSubheading
        {Özel ALEV Okulları}{September 2008 -- Juni 2017}
        {Grund- und Mittelschule}{Istanbul, Türkei}
        \resumeItemListStart
       \small\resumeItem{TEOG (landesweite Prüfung): unter den besten 0,007\,\% der Türkei}
        \resumeItemListEnd
\resumeSubHeadingListEnd


%-----------ERFAHRUNG-----------
\section{Erfahrung}
\resumeSubHeadingListStart
\resumeSubheading
        {Sensorfusion zur Bewertung von Patientenbewegungen}{Juni 2026 -- August 2026}
        {TUM Lehrstuhl für Kognitive Systeme - Studentische Hilfskraft}{München, Deutschland}
        \vspace{0.09cm}
        \resumeItemListStart
        \small\resumeItem{Entwicklung von Sensorfusionsmethoden für tragbare IMU-Körpersensoren zur quantitativen Bewertung der Verbesserung oder Verschlechterung der Motorik von Patienten.}
        \vspace{0.05cm}
        \resumeItem{Konzeption der Kommunikationsarchitektur zwischen verteilten ESP32-Mikrocontrollern für Echtzeit-Synchronisation und Datenaustausch in einem skalierbaren Netzwerk von ca. 20 MCUs.}
        \resumeItemListEnd
          \vspace{0.2cm}
\resumeSubheading
        {Embedded AI in Echtzeitanwendungen}{November 2024 -- April 2026}
        {Infineon Technologies - Werkstudent}{München, Deutschland}
        \vspace{0.09cm}
        \resumeItemListStart
        \small\resumeItem{Entwicklung von Echtzeit-KI-Systemen für sensorlose Motorregelung auf Edge-Geräten, u.\,a. Echtzeit-Anomalieerkennung für die Motorsicherheit und sensorlose Positionsschätzung auf dem Chip.}
        \vspace{0.05cm}
        \resumeItem{Mitarbeit an der Entwicklung und Portierung von Umgebungslösungen für AURIX-Mikrocontroller.}
        \resumeItemListEnd
          \vspace{0.2cm}
\resumeSubheading
        {NEURA Robotics Challenge}{April 2025 -- Oktober 2025}
        {BCI-gesteuertes Robotikprojekt - Teamleiter}{München, Deutschland}
        \vspace{0.09cm}
        \resumeItemListStart
        \small\resumeItem{Teilnahme an der Forschungs-Challenge von NEURA Robotics, bei der jedes Team eine Roboterhand für eine kreative Aufgabe erhielt.}
        \vspace{0.05cm}
        \resumeItem{Entwicklung einer SSVEP-basierten Echtzeit-Pipeline für ein Brain-Computer-Interface (BCI) auf ROS zur Steuerung eines MAIRA-Roboterarms.}
        \vspace{0.05cm}
        \resumeItem{Einer von 6 Finalisten; mit 10.000\,€ Preisgeld ausgezeichnet, das nun die Weiterentwicklung des Projekts finanziert.}
        \resumeItemListEnd
          \vspace{0.2cm}
\resumeSubheading
        {Tutor - Systemtheorie}{April 2024 -- August 2024}
        {Methoden der Signalverarbeitung - Dr.-Ing. Michael Joham}{München, Deutschland}
        \vspace{0.09cm}
        \resumeItemListStart
        \small\resumeItem{Initiierung einer Vorlesungsreihe mit über 80 Teilnehmenden zur Vertiefung des Verständnisses der Systemtheorie.}
        \vspace{0.05cm}
        \resumeItem{Vermittlung effektiver Lerngewohnheiten und Prüfungsstrategien, wodurch die Bestehensquote um 30\,\% stieg.}
        \resumeItemListEnd
          \vspace{0.2cm}
\resumeSubheading
        {Advisor-Programm (Robotics Race)}{Oktober 2022 -- Juli 2023}
        {Teamleiter}{München, Deutschland}
        \vspace{0.09cm}
        \resumeItemListStart
        \small\resumeItem{Leitung eines 10-köpfigen Teams beim Bau eines Küchenroboters, der individuelle Sandwiches zubereitet, inkl. Moderation von Meetings und Festlegung von Meilensteinen.}
        \vspace{0.05cm}
        \resumeItem{Konstruktion der Gehäusekomponenten per 3D-Modellierung und Programmierung des Mikrocontrollers in C.}
        \resumeItemListEnd
\resumeSubHeadingListEnd


%-----------AKTIVITÄTEN & ENGAGEMENT-----------
\section{Aktivitäten und Engagement}
\resumeSubHeadingListStart
\resumeSubheading
        {NeuroTUM - Neurowissenschaftlicher Verein an der TUM}{Oktober 2024 -- heute}
        {Co-Director \& Teamleiter von neuromotion}{München, Deutschland}
      \resumeItemListStart
      \vspace{0.1cm}
        \small\resumeItem{Mitgründung und Leitung von neuromotion: Entwicklung eines assistiven Roboterarms, gesteuert über ein Echtzeit-BCI, damit tetraplegische Menschen Alltagsaufgaben selbstständig bewältigen können.}
        \vspace{0.05cm}
        \resumeItem{Als Co-Director Aufbau der Neuroscience-Community an der TUM durch Koordination der Teamkommunikation, Einführung neuer Tools und Arbeitsweisen sowie Öffentlichkeitsarbeit.}
        \vspace{0.05cm}
        \resumeItem{Vertretung von NeuroTUM auf Messen und Veranstaltungen mit Vorträgen über unsere Projekte und Forschung.}
    \resumeItemListEnd
    \vspace{0.1cm}
\resumeSubheading
        {Munich Neuromorphic Hackathon}{November 2025 \& Oktober 2026}
        {Organisator der 3. \& 4. Ausgabe, gemeinsam mit fortiss}{München, Deutschland}
      \resumeItemListStart
      \vspace{0.1cm}
        \small\resumeItem{Heranführung interessierter Studierender an Neuromorphic Computing und neue Rechenkonzepte.}
    \resumeItemListEnd
\resumeSubHeadingListEnd


%-----------LEHRVERANSTALTUNGEN & PROJEKTE-----------
\section{Relevante Lehrveranstaltungen \& Projekte}
\resumeSubHeadingListStart
\resumeSubheading
        {Assessment of Neurological Diseases}{Januar 2024 -- Juni 2024}
        {TUM Projektwoche}{München, Deutschland}
        \resumeItemListStart
        \small\resumeItem{Entwicklung einer C\#-App zur Datenerfassung bei Patienten mit Multipler Sklerose.}
        \resumeItemListEnd
        \vspace{0.05cm}
\item \small\textbf{Ausgewählte Module}: Schaltungstheorie, Lineare Algebra, Blockpraktikum C++, Praktikum VHDL, Analysis 1--3, Systeme der Signalverarbeitung, Introduction to Deep Learning, Echtzeitsysteme, Machine Learning, 3D Computer Vision, Robot Dynamics
\resumeSubHeadingListEnd


%-----------ZERTIFIKATE & AUSZEICHNUNGEN-----------
\section{Zertifikate und Auszeichnungen}
\resumeSubHeadingListStart
\resumeSubheading
        {AI Programming with Python - Bertelsmann}{Dezember 2023 -- Juni 2024}
        {Stipendiat der Online-Plattform Udacity}{Online, Selbststudium}\break
        \resumeItemListStart
         \small\resumeItem{{\url{https://t.ly/5EFDU}}}
        \resumeItemListEnd
        \vspace{0.1cm}
\resumeSubheading
        {Erster Platz im Advisor-Wettbewerb}{Juli 2023}
        {Technische Universität München}{München, Deutschland}\break
      \vspace{0.1cm}
\resumeSubheading
        {Object-oriented Programming}{Oktober 2022}
        {Technische Universität München}{Online, Selbststudium}
      \resumeItemListStart
        \small\resumeItem{{\url{https://t.ly/N-2HW}}}
    \resumeItemListEnd
\resumeSubheading
        {Goethe-Zertifikat C1}{Dezember 2021}
        {Goethe-Institut}{Istanbul, Türkei}
\resumeSubHeadingListEnd


%-----------KENNTNISSE & SPRACHEN-----------
\section{Kenntnisse \& Sprachen}
\small
\begin{tabular*}{\textwidth}{@{\extracolsep{\fill}} l}
\textbf{Programmierung}: C/C++, Python (PyTorch, Pandas, Matplotlib, TensorFlow, pylsl), Java,\\ MATLAB, C\#, HTML/CSS \\
\vspace{0.1cm}
\textbf{Embedded}: ESP32 (ESP-IDF), AURIX Development Studio, Arduino, Löten \\
\vspace{0.1cm}
\textbf{Robotik}: ROS2, RViz, Isaac Sim \\
\vspace{0.1cm}
\textbf{Tools}: Git/GitHub, VS Code, Fusion, Qt Creator, Anaconda, MS Office \\
\vspace{0.1cm}
\textbf{Betriebssysteme}: Linux, Windows, macOS \\
\vspace{0.1cm}
\textbf{Sprachen}: Türkisch (C2 / Muttersprache), Deutsch (C1 / Fließend), Englisch (C1 / Fließend)
\end{tabular*}

\end{document}
